\section{SmallWood Programming Model}
\label{sec:smallwood-model}
\label{sec:smallwood}

This section specifies the SmallWood relations used by ARC-SPRUCE. Its purpose is to record the committed witness
surfaces, public inputs, and PACS constraint families used by the issuance and showing proofs, and to state the precise
strength of the resulting proof statements. Throughout, SmallWood certifies only the compiled relations stated in this
section, namely relations over committed witness rows, public lifts, and PACS constraint families derived from them.
Any further coefficient-domain consequence is obtained only through the explicit theorem stated below.

\subsection{Proof-language interface}
\label{subsec:smallwood-interface}
Fix an evaluation set \(\Omega=\{\omega_1,\ldots,\omega_s\}\subset \Fq\) and a disjoint set of blinding points
\(\Omega'\subset \Fq\setminus \Omega\). A witness is arranged as a matrix
\[
W\in \Fq^{n\times s},
\]
whose rows are interpolated into low-degree polynomials
\[
P_1(X),\ldots,P_n(X)\in \Fq[X]
\]
satisfying
\[
P_i(\omega_k)=W_{i,k}\qquad\text{for all }\omega_k\in\Omega,
\]
and taking random values on \(\Omega'\) in order to achieve honest-verifier zero knowledge. SmallWood commits to these
row polynomials and proves satisfaction of PACS constraints by opening verifier-chosen linear maps of the committed rows.

The PACS language distinguishes two classes of constraints.
\begin{itemize}[leftmargin=1.25em]
    \item \emph{Parallel constraints} are enforced pointwise on \(\Omega\). They are used for slot-local algebra,
    including linear relations, pointwise products, selector gating, and checkpointed PRF constraints.
    \item \emph{Aggregated constraints} are enforced through \(\Omega\)-sum identities. They are used when the encoding
    introduces global coupling across packed coordinates, in particular for public linear bridge relations between source
    rows and replay-image rows.
\end{itemize}

The proof-system properties imported from SmallWood apply to the compiled PACS relation built from these rows and
constraints. They do not, by themselves, replace the explicit coefficient-domain consequence theorem stated in
Section~\ref{subsec:statement-strength}.

\subsection{Abstract and compiled relations}
\label{subsec:abstract-compiled}
At the protocol layer, the blind-signature and ARC constructions are specified by abstract relations in the quotient ring
\(R_q\). The SmallWood proofs, however, are proofs for compiled relations over committed witness rows, public lifts, and
PACS constraints. We therefore keep the two levels distinct throughout this section.

In issuance, the compiled relation is the cleared-denominator relation
\(R_{\mathrm{Issue}}^{\mathrm{cmp}}\) of Section~\ref{sec:blind-signature}, with the public target \(T\) supplied as an
input to the proof. In showing, the compiled relation is a full transform-basis replay relation over explicitly committed
source rows, full replay-image blocks, and PRF companion rows. The purpose of the remainder of the section is to state
those compiled relations precisely and to prove their coefficient-domain consequence.

\subsection{Issuance (pre-sign) compiled relation}
\label{subsec:smallwood-issuance}
\label{subsec:presign-compiled}
The pre-sign proof keeps the blind-signature target \(T\) public and proves that it is consistent with one hidden
commitment opening and the issuer randomness. At the abstract level, issuance is governed by
\(R_{\mathrm{Issue}}^{\mathrm{abs}}\) from Section~\ref{sec:blind-signature}. The compiled SmallWood proof certifies
instead the weaker relation \(R_{\mathrm{Issue}}^{\mathrm{cmp}}\), which replaces the defined quotient semantics by the
cleared-denominator relation.

\paragraph{Committed rows.}
The committed pre-sign witness consists of the following rows:
\[
C^M,\ C^{\mathrm{preRU}},\ C^{\mathrm{preR}},\ C^{\mathrm{ctr}},\ C^J,
\]
\[
M,\ K,\ R_0^U,\ R_1^U,\ R,\ R_0,\ R_1,\ J_0,\ J_1,
\]
\[
\widehat M,\ \widehat K,\ \widehat R_0,\ \widehat R_1.
\]
The first five rows are low-alphabet carriers. The next nine rows are decoded aliases. The final four rows are the
transform aliases used by the hash replay.

\paragraph{Public inputs and lifts.}
The public pre-sign inputs are
\[
\mACom,\ B,\ \vecc,\ \vecr_0^{\Issuer},\ \vecr_1^{\Issuer},\ T.
\]
The verifier also derives the public lifts
\[
B_0^\Theta,\ B_1^\Theta,\ B_2^\Theta,\ B_3^\Theta,\ T^\Theta,
\]
where \(T^\Theta\) is the chosen public lift of the coefficient-domain target \(T\).

\paragraph{Constraint families.}
Let \(\Delta=2\BoundMsg+1\). The compiled pre-sign relation enforces:
\begin{enumerate}[leftmargin=1.25em]
    \item \emph{Carrier membership:}
    \[
    P_M(C^M)=0,\qquad
    P_{\mathrm{preRU}}(C^{\mathrm{preRU}})=0,\qquad
    P_{\mathrm{preR}}(C^{\mathrm{preR}})=0,
    \]
    \[
    P_{\mathrm{ctr}}(C^{\mathrm{ctr}})=0,\qquad
    P_J(C^J)=0.
    \]
    \item \emph{Carrier-to-alias bridges:}
    \[
    M-\mathrm{Dec}_M^{(1)}(C^M)=0,\qquad
    K-\mathrm{Dec}_M^{(2)}(C^M)=0,
    \]
    \[
    R_0^U-\mathrm{Dec}_{\mathrm{preRU}}^{(1)}(C^{\mathrm{preRU}})=0,\qquad
    R_1^U-\mathrm{Dec}_{\mathrm{preRU}}^{(2)}(C^{\mathrm{preRU}})=0,
    \]
    \[
    R-\mathrm{Dec}_{\mathrm{preR}}^{(1)}(C^{\mathrm{preR}})=0,
    \]
    \[
    R_0-\mathrm{Dec}_{\mathrm{ctr}}^{(1)}(C^{\mathrm{ctr}})=0,\qquad
    R_1-\mathrm{Dec}_{\mathrm{ctr}}^{(2)}(C^{\mathrm{ctr}})=0,
    \]
    \[
    J_0-\mathrm{Dec}_J^{(1)}(C^J)=0,\qquad
    J_1-\mathrm{Dec}_J^{(2)}(C^J)=0.
    \]
    \item \emph{Commitment relation:}
    \[
    \mACom\cdot [M\Vert K\Vert R_0^U\Vert R_1^U\Vert R]^\top = \vecc.
    \]
    \item \emph{Centering relation:}
    \[
    R_0 = R_0^U + \vecr_0^{\Issuer} - \Delta J_0,\qquad
    R_1 = R_1^U + \vecr_1^{\Issuer} - \Delta J_1.
    \]
    \item \emph{Packing relation:} the rows \(M\) and \(K\) occupy complementary halves of the packing domain and
    together encode the abstract message \(m=\mu\Vert k\).
    \item \emph{Transform-domain cleared relation:}
    \[
    (B_3^\Theta-\widehat R_1)\Had T^\Theta
    -
    \bigl(B_0^\Theta + B_1^\Theta\Had(\widehat M+\widehat K)+B_2^\Theta \Had \widehat R_0\bigr)
    =0.
    \]
\end{enumerate}

\paragraph{Aggregated bridge families.}
The issuance statement includes the four transform bridges
\[
M\mapsto \widehat M,\qquad
K\mapsto \widehat K,\qquad
R_0\mapsto \widehat R_0,\qquad
R_1\mapsto \widehat R_1.
\]

\paragraph{Statement proved.}
The compiled pre-sign proof asserts the existence of committed carriers, decoded aliases, and transform aliases such
that:
\begin{enumerate}[leftmargin=1.25em]
    \item each carrier lies in its designated carrier alphabet;
    \item each decoded alias is the decode image of its carrier;
    \item the decoded aliases satisfy commitment, centering, and packing;
    \item the transform aliases are the public transform images of the decoded rows; and
    \item the public target \(T\), after lifting, satisfies the transform-domain cleared relation.
\end{enumerate}
This is a proof of \(R_{\mathrm{Issue}}^{\mathrm{cmp}}\). No inverse witness for \(B_3-R_1\) is included, and the
compiled statement should therefore not be described as a proof of the abstract quotient semantics.

\subsection{Showing (post-sign) compiled relation}
\label{subsec:smallwood-showing}
\label{subsec:showing-compiled}
At the ARC layer, the abstract showing relation is the existence of a witness \((u,\mu,k,r_0,r_1)\) such that
\[
Au = h_{\mu\Vert k,(r_0,r_1)}(B)
\qquad\text{and}\qquad
\prftag = \PRF(k,\nonce),
\]
together with the public coefficient-wise bounds on \(\mu,k,r_0,r_1\) and the public shortness bound on \(u\). The
post-sign NIZK does not expose this witness directly. Instead, it proves a compiled transform-basis relation over
committed source rows, full replay-image blocks, and PRF companion rows.

\paragraph{Exact replay transform.}
Let
\[
\mathcal F_{\mathrm{rep}} : \Rq \longrightarrow \mathcal T
\]
denote the exact public transform used by the replay relation, where \(\mathcal T\) is the full transform image of
\(\Rq\) for the chosen replay basis. We write
\[
\mathcal F_{\mathrm{rep}} = \bigl(\mathcal F_{\mathrm{rep},0},\ldots,\mathcal F_{\mathrm{rep},B_{\mathrm{rep}}-1}\bigr)
\]
for its decomposition into replay blocks. For every public hash polynomial \(B_i\in\Rq\), define the corresponding
block lifts
\[
B_{i,b}^\Theta := \mathcal F_{\mathrm{rep},b}(B_i),
\qquad i\in\{0,1,2,3\},\quad b\in\{0,\ldots,B_{\mathrm{rep}}-1\}.
\]

\paragraph{Committed rows.}
The committed showing witness consists of:
\begin{enumerate}[leftmargin=1.25em]
    \item the signature limb rows \(L_{c,b,\ell}\);
    \item the carrier rows \(C^M\) and \(C^{\mathrm{ctr}}\);
    \item the decoded non-sign source rows
    \[
    M,\ K,\ R_0,\ R_1;
    \]
    \item the full replay-image block families
    \[
    \widehat T_0,\ldots,\widehat T_{B_{\mathrm{rep}}-1},
    \qquad
    \widehat m_0,\ldots,\widehat m_{B_{\mathrm{rep}}-1},
    \]
    \[
    \widehat R_{0,0},\ldots,\widehat R_{0,B_{\mathrm{rep}}-1},
    \qquad
    \widehat R_{1,0},\ldots,\widehat R_{1,B_{\mathrm{rep}}-1};
    \]
    \item the PRF companion rows.
\end{enumerate}
The coefficient-domain signature witness \(u\) is reconstructed from the limb surface by the public digit-collapse map
of the shortness gadget.

\paragraph{Parallel families.}
The compiled showing relation enforces:
\begin{enumerate}[leftmargin=1.25em]
    \item \emph{Signature digit membership:} the limb rows satisfy the signed low-alphabet digit constraints of the
    chosen shortness profile.
    \item \emph{Carrier membership:}
    \[
    P_M(C^M)=0,\qquad P_{\mathrm{ctr}}(C^{\mathrm{ctr}})=0.
    \]
    \item \emph{Carrier-to-source decode:}
    \[
    M=\mathrm{Dec}_M^{(1)}(C^M),\qquad K=\mathrm{Dec}_M^{(2)}(C^M),
    \]
    \[
    R_0=\mathrm{Dec}_{\mathrm{ctr}}^{(1)}(C^{\mathrm{ctr}}),\qquad
    R_1=\mathrm{Dec}_{\mathrm{ctr}}^{(2)}(C^{\mathrm{ctr}}).
    \]
    \item \emph{Packing:} the decoded rows \(M\) and \(K\) satisfy the same half-split message representation as in
    issuance.
    \item \emph{Full replay residual:} for every block
    \[
    b\in\{0,\ldots,B_{\mathrm{rep}}-1\},
    \]
    the replay aliases satisfy
    \[
    (B_{3,b}^\Theta-\widehat R_{1,b})\Had \widehat T_b
    -
    \bigl(B_{0,b}^\Theta + B_{1,b}^\Theta\Had \widehat m_b + B_{2,b}^\Theta\Had \widehat R_{0,b}\bigr)
    =0.
    \]
    \item \emph{PRF key binding and nonlinear constraints:} the companion rows are bound to the decoded
    \(K\)-coordinates, satisfy the checkpointed permutation equations, and bind the public tag after feed-forward and
    truncation.
\end{enumerate}

\paragraph{Aggregated bridge families.}
The compiled showing relation includes the following aggregated families.
\begin{enumerate}[leftmargin=1.25em]
    \item \emph{Signature-side bridges:} for every replay block \(b\),
    \[
    \widehat T_b = \mathcal F_{\mathrm{rep},b}(Au),
    \]
    where \(u\) is the coefficient-domain signature witness reconstructed from the cert-side limb surface.
    \item \emph{Non-sign transform bridges:} for every replay block \(b\),
    \[
    \widehat m_b = \mathcal F_{\mathrm{rep},b}(M+K),\qquad
    \widehat R_{0,b} = \mathcal F_{\mathrm{rep},b}(R_0),\qquad
    \widehat R_{1,b} = \mathcal F_{\mathrm{rep},b}(R_1).
    \]
\end{enumerate}

\paragraph{Statement proved.}
The compiled showing proof asserts the existence of committed limb rows, source rows, replay-image rows, and PRF
companion rows such that:
\begin{enumerate}[leftmargin=1.25em]
    \item the cert-side signature witness is short via the limb decomposition;
    \item the decoded source rows are the coefficient-domain message and randomness rows used throughout the statement;
    \item the full replay-image blocks are the exact transform images of \(Au\), \(M+K\), \(R_0\), and \(R_1\);
    \item the full replay residual vanishes blockwise; and
    \item the public tag is bound to the signed hidden key and the public nonce by the PRF companion constraints.
\end{enumerate}

\subsection{PRF constraints}
\label{subsec:prf-checkpointed}
\label{subsec:prf-constraints}
The public tag is tied to the signed hidden key by a checkpointed arithmetization of the Poseidon2-like permutation from
Section~\ref{subsec:poseidon-prf}. The proof commits selected S-box outputs and lets the verifier replay all public
linear wiring between checkpoints from openings and public constants. The resulting families are parallel PACS
constraints over the companion rows, together with the public feed-forward and truncation equations that bind the final
state to the tag.

\subsection{Coefficient-domain consequence and statement strength}
\label{subsec:statement-strength}
The completeness, knowledge-soundness, and zero-knowledge guarantees imported from SmallWood apply to the compiled PACS
relations stated in this section. In particular, the showing proof directly certifies the full transform-basis replay
relation on explicitly committed replay-image blocks together with the signature-side and non-sign bridge families. The
coefficient-domain consequence is the following.

\begin{lemma}[Exact replay transform]\label{lem:replay-transform}
The public replay map
\[
\mathcal F_{\mathrm{rep}} : \Rq \longrightarrow \mathcal T
\]
is the exact transform isomorphism used by the replay basis. In particular, it is additive, multiplicative, and
injective.
\end{lemma}

\begin{proof}
The replay basis is chosen so that \(\mathcal F_{\mathrm{rep}}\) is the corresponding CRT/NTT isomorphism of
\(\Rq\). The stated properties are exactly the defining properties of that isomorphism.
\end{proof}

\begin{lemma}[Coefficient-domain consequence of the compiled showing relation]\label{lem:showing-to-cleared}
Assume the compiled showing relation of Section~\ref{subsec:showing-compiled}. Then
\[
B_3Au-(Au)R_1-B_0-B_1(M+K)-B_2R_0=0
\]
in \(\Rq\).
\end{lemma}

\begin{proof}
For every replay block \(b\), the bridge families and the public lift definition give
\[
\widehat T_b = \mathcal F_{\mathrm{rep},b}(Au),\qquad
\widehat m_b = \mathcal F_{\mathrm{rep},b}(M+K),
\]
\[
\widehat R_{0,b} = \mathcal F_{\mathrm{rep},b}(R_0),\qquad
\widehat R_{1,b} = \mathcal F_{\mathrm{rep},b}(R_1),
\]
\[
B_{i,b}^\Theta = \mathcal F_{\mathrm{rep},b}(B_i)\qquad\text{for }i\in\{0,1,2,3\}.
\]
Substituting these identities into the blockwise replay residual gives
\[
\mathcal F_{\mathrm{rep},b}(B_3Au)
-
\mathcal F_{\mathrm{rep},b}((Au)R_1)
-
\mathcal F_{\mathrm{rep},b}(B_0)
-
\mathcal F_{\mathrm{rep},b}(B_1(M+K))
-
\mathcal F_{\mathrm{rep},b}(B_2R_0)
=0
\]
for every \(b\). Concatenating the blocks yields
\[
\mathcal F_{\mathrm{rep}}
\bigl(
B_3Au-(Au)R_1-B_0-B_1(M+K)-B_2R_0
\bigr)=0.
\]
By Lemma~\ref{lem:replay-transform}, \(\mathcal F_{\mathrm{rep}}\) is injective. Hence
\[
B_3Au-(Au)R_1-B_0-B_1(M+K)-B_2R_0=0
\]
in \(\Rq\).
\end{proof}

\begin{corollary}[Abstract rational-hash relation under admissibility]\label{cor:showing-to-abstract}
If, in addition, the witness satisfies
\[
B_3-R_1\in \Rq^\times,
\]
then
\[
Au = h_{M+K,(R_0,R_1)}(B).
\]
\end{corollary}

\begin{proof}
Lemma~\ref{lem:showing-to-cleared} gives
\[
(B_3-R_1)\Had Au - \bigl(B_0+B_1(M+K)+B_2R_0\bigr)=0.
\]
Under the admissibility condition \(B_3-R_1\in \Rq^\times\), Lemma~\ref{lem:hash-cleared-admissible} implies
\[
Au = h_{M+K,(R_0,R_1)}(B).
\]
\end{proof}

\begin{corollary}[Consequence of an accepting post-sign proof]\label{cor:showing-proof-consequence}
Assume the post-sign NIZK is knowledge sound for the compiled showing relation. Then every accepting post-sign proof
yields a witness \((u,M,K,R_0,R_1)\) satisfying
\[
B_3Au-(Au)R_1-B_0-B_1(M+K)-B_2R_0=0
\]
in \(\Rq\). If the extracted witness additionally satisfies denominator admissibility, then it also satisfies
\[
Au = h_{M+K,(R_0,R_1)}(B).
\]
\end{corollary}

\begin{proof}
Knowledge soundness yields a witness satisfying the compiled showing relation. Lemma~\ref{lem:showing-to-cleared}
then gives the coefficient-domain cleared identity, and Corollary~\ref{cor:showing-to-abstract} gives the quotient-side
relation under admissibility.
\end{proof}

\begin{remark}[Statement strength]
The theorem suite above proves a one-way consequence from the compiled full-replay showing relation to the
coefficient-domain cleared identity, and from there to the abstract quotient relation under admissibility. It does not
assert a converse implication from the coefficient-domain identity back to the full compiled relation. The present paper
therefore does not claim a two-way equivalence between the compiled showing relation and the abstract showing relation.
\end{remark}

\begin{remark}[Strengthened post-sign variant]\label{rem:showing-admissibility-variant}
If one wishes the post-sign proof itself to certify the abstract quotient relation without any external admissibility
side condition, the compiled statement can be strengthened by introducing an inverse witness \(z\) and enforcing
\[
(B_3-R_1)\Had z = 1,\qquad
Au = z\Had \bigl(B_0+B_1(M+K)+B_2R_0\bigr).
\]
The present section does not assume that this stronger variant has been adopted.
\end{remark}

The remaining sections use Lemma~\ref{lem:showing-to-cleared} and Corollary~\ref{cor:showing-proof-consequence}
whenever a coefficient-domain consequence of the compiled showing proof is needed, and invoke
Corollary~\ref{cor:showing-to-abstract} only when admissibility is also available.
